\section{Metodologia de la practica}

\subsection{Actividad 1: Cambio de color de fondo}

Para resolver el primer ejercicio, se requirio la conversión del formato de color. A diferencia de las herramientas de diseño convencionales, OpenGL no interpreta los colores en hexadecimal ni en la escala RGB de 0 a 255, sino que requiere valores de punto flotante normalizados en un rango de 0.0f a 1.0f.

El algoritmo de conversión matemática que aplicamos para la práctica consistió en dos pasos elementales:
\begin{itemize}
    \item Separar el código hexadecimal de la paleta en sus componentes independientes (rojo, verde y azul) y transformarlos a base decimal.
    \item Dividir cada uno de estos valores decimales entre 255.0 para obtener su equivalente normalizado como un número flotante.
\end{itemize}

Tomando como ejemplo el primer color solicitado, \textbf{\#00C2CB}, el canal rojo (00) se mantiene en 0.0f. El canal verde (C2) equivale a 194 en decimal, que dividido entre 255 nos da 0.7607f. Finalmente, el canal azul (CB) es 203, arrojando un valor de 0.7960f. Iteramos esta misma rutina matemática para extraer los valores de los cinco tonos de la paleta exigida.

A nivel de arquitectura de software, trabajamos sobre el archivo principal \texttt{01\_IntroOpenGL.cpp}. Para que el color reaccionara en tiempo de ejecución, nos enganchamos a la función \texttt{keyboardCallback} del código base, la cual gestiona las interrupciones de hardware.

Tal como lo indica la rúbrica de evaluación[cite: 1], a continuación se presenta el segmento de código relevante con la implementación. Mapeamos las teclas R, G, B, Y, y M para sobreescribir dinámicamente las variables globales \texttt{bgR}, \texttt{bgG}, \texttt{bgB} y \texttt{bgA}. En cada ciclo de fotograma, la función \texttt{Update()} lee estas variables y se las pasa al comando nativo \texttt{glClearColor}, logrando así el barrido y pintado de la pantalla con el color seleccionado:\\

Código de los cambios realizados para el cambio de color de fondo:\\

\begin{lstlisting}[language=C, caption=Código del ejercicio 1]
        // Actividad 1.1
    if (glfwGetKey(window, GLFW_KEY_R) == GLFW_PRESS) {
        bgR = 0.0f;
        bgG = 0.7607f;
        bgB = 0.7960f;
        bgA = 1.0;
    }
    if (glfwGetKey(window, GLFW_KEY_G) == GLFW_PRESS) {
        bgR = 0.4980f;
        bgG = 0.8784f;
        bgB = 0.8313f;
        bgA = 1.0;
    }
    if (glfwGetKey(window, GLFW_KEY_B) == GLFW_PRESS) {
        bgR = 1.0f;
        bgG = 0.8901f;
        bgB = 0.5411f;
        bgA = 1.0;
    }
    if (glfwGetKey(window, GLFW_KEY_N) == GLFW_PRESS) {
        bgR = 1.0f;
        bgG = 0.6039f;
        bgB = 0.4627f;
        bgA = 1.0;
    }
    if (glfwGetKey(window, GLFW_KEY_O) == GLFW_PRESS) {
        bgR = 1.0f;
        bgG = 0.3647f;
        bgB = 0.5607f;
        bgA = 1.0;
    }
\end{lstlisting}
